% =============================================================
%  ch19_kamera_anbindung.tex  --  Kamera anbinden, Frame splitten
% =============================================================

\chapter{Kamera anbinden, Frame splitten}

Die synchronisierte ELP-Kamera liefert ein side-by-side-Bild. Erst den Treiber
starten, dann links/rechts trennen.
\begin{lstlisting}[style=bash,caption={Treiber starten und Topics prüfen}]
ros2 run v4l2_camera v4l2_camera_node \
  --ros-args -p video_device:=/dev/video0 \
             -p image_size:="[2560,720]"
ros2 topic list                          # /image_raw sollte erscheinen
ros2 run rqt_image_view rqt_image_view   # Bild ansehen
\end{lstlisting}

\begin{lstlisting}[style=py,caption={Knoten: side-by-side Frame in /left und /right splitten}]
import rclpy
from rclpy.node import Node
from sensor_msgs.msg import Image
from cv_bridge import CvBridge

class StereoSplitter(Node):
    def __init__(self):
        super().__init__('stereo_splitter')
        self.bridge = CvBridge()
        self.pub_l = self.create_publisher(Image, '/left/image_raw', 10)
        self.pub_r = self.create_publisher(Image, '/right/image_raw', 10)
        self.create_subscription(Image, '/image_raw', self.cb, 10)

    def cb(self, msg):
        frame = self.bridge.imgmsg_to_cv2(msg, 'bgr8')
        w = frame.shape[1] // 2
        left, right = frame[:, :w], frame[:, w:]
        hdr = msg.header  # keep timestamp so stereo stays synchronized
        lm = self.bridge.cv2_to_imgmsg(left,  'bgr8'); lm.header = hdr
        rm = self.bridge.cv2_to_imgmsg(right, 'bgr8'); rm.header = hdr
        self.pub_l.publish(lm)
        self.pub_r.publish(rm)

def main():
    rclpy.init()
    rclpy.spin(StereoSplitter())
    rclpy.shutdown()
\end{lstlisting}
